%!TEX root=../../main.tex

\begin{restatable}{lemma}{boundedSolutionsCGL}\label{lemma:bounded-solutions-cgl}
	Every solution to the constrained group lasso problem
	is bounded by an absolute constant independent of \( \lambda \).
	Specifically, every \( w \in \solfn(\lambda) \) satisfies
	\[
		\sum_{\bi \in \calB} \norm{\wi}_2 \leq \sum_{\bi \in \calB} \norm{\bar w_\bi}_2.
	\]
	where \( \bar w \in \solfn(0) \) is the least-squares solution with minimum \( \ell_2 \)-norm.
\end{restatable}
\begin{proof}
	Let \( h(w) = \sum_{\bi \in \calB} \norm{\wi}_2 \),
	and define \( W_{g} \) to be the set of least squares solutions with
	minimum group norm,
	\[
		W_{g} = \argmin \cbr{h(w) : w \in \solfn(0) }.
	\]
	Let \( w_g \in W_g \),
	and suppose that \( h(w_{g}) < h(w(\lambda)) \)
	for some \( \lambda > 0 \), \( w \in \solfn(\lambda) \).
	Since
	\[
		\half \norm{X w_g - y}_2^2
		\leq \half \norm{X w(\lambda) - y}_2^2
	\]
	we deduce
	\[
		\half \norm{X w_g - y}_2^2 + \lambda h(w_g) <
		\half \norm{X \w(\lambda) - y}_2^2
		+ \lambda h(\w(\lambda)),
	\]
	which is a contradiction.
	So \( h(\w(\lambda)) \leq h(w_g) \) for all \( \lambda > 0 \).
	Observing \( h(w_g) \leq h(\bar w) \) since \( \bar w \) may not be in
	\( W_{g} \) gives the result.
	Since \( h(\bar w) \) is independent of \( \lambda \),
	we conclude that \( \solfn(\lambda) \) is bounded independent
	of \( \lambda \).
\end{proof}

\valueContinuityCGL*
\begin{proof}
	Define the joint objective function
	\[
		f(w, \lambda) = \half \norm{X w - y}_2^2 + \lambda \sum_{b_i \in \calB} \norm{\wi}_2.
	\]
	Clearly \( f(w, \lambda) \) is jointly continuous in \( w \) and \( \lambda \).
	By \cref{lemma:bounded-solutions-cgl}, minimization of
	\( f(w, \lambda) \) subject to \( \Ki^\top \wi \leq 0 \)
	is equivalent to the constrained minimization problem,
	\[
		p^*(\lambda) = \min_{w} f(w, \lambda) \quad \text{s.t.} \; \sum_{\bi \in \calB} \norm{\wi}_2 \leq C, \, \Ki^\top \wi \leq 0,
	\]
	where \( C \) is a finite absolute constant.
	Note that this expression is also valid when \( \lambda = 0 \) as the
	min-norm solution to the unregularized least squares problem
	obeys the constraint.

	Thus, \( \solfn \) is a continuous optimization problem over a continuous
	(constant in this case)
	compact constraint set and the classical result of \citet{berge1997topological}
	(see also \citet{hogan1973point}[Theorem 7]) implies
	\( p^* \) is continuous.
\end{proof}

\mapContinuityCGL*
\begin{proof}
	Joint continuity of the objective
	\[
		f(w, \lambda) = \half \norm{X w - y}_2^2 + \lambda \sum_{b_i \in \calB} \norm{\wi}_2,
	\]
	combined with continuity of the (constant) constraint allows us to use
	\citet[Theorem 1]{robinson1974sufficient} to obtain that
	that \( \solfn \) is upper semi-continuous.
	Since \( \solfn \) is convex and bounded, it is compact.
	It is thus also uniformly compact and upper semi-continuity is
	equivalent to closedness \citep{hogan1973point}[Theorem 3].
	We conclude that \( \solfn \) is closed as claimed.

	If \( X \) is full column rank, then the constrained group
	lasso solution is unique for all \( \lambda \geq 0 \).
	The solution map is a singleton on \( \R^+ \), and closedness
	and openness are equivalent properties for singleton maps.
	Since we have already shown it is closed, the solution map
	must also be open.
	An identical argument shows that the solution map is open
	at on any interval over which the solution is unique.

	Now we show the reverse implication by proving the contrapositive.
	Assume \( X \) is not full column-rank and suppose \( \Ki = 0 \)
	for each \( \bi \in \calB \).
	The solution map at \( \lambda = 0 \) is the solution set to
	the least squares problem,
	\[
		\min_w \half \norm{X w - y}_2^2,
	\]
	which is known to be \( \solfn(0) = \cbr{\wmin(0) + z : z \in \Null(X)} \).
	While \( \solfn(0) \) is unbounded, it holds that
	\( \solfn(\lambda) \subset C \) for some bounded \( C \) for
	every \( \lambda > 0 \) (\cref{lemma:bounded-solutions-cgl}).
	As a result, there exist uncountably many solutions in \( \solfn(0) \)
	which are not limit points of solutions in \( \solfn(\lambda_k) \)
	as \( \lambda_k \rightarrow 0 \).
	In other words, \( \solfn(0) \) is not open at \( 0 \).
\end{proof}


\modelFitContinuity*
\begin{proof}
	Consider the dual problem from \cref{lemma:lagrange-dual},
	\begin{equation*}
		\begin{aligned}
			\max_{\eta} \;
			 & - \half (\eta - X^\top y) (X^\top X)^+ (\eta - X^\top y) + \half \norm{y}_2^2                                  \\
			 & \quad \quad \text{s.t.} \quad \eta \in \Row(X), \; \norm{\eta_\bi}_2 \leq \lambda \; \forall \; \bi \in \calB.
		\end{aligned}
	\end{equation*}
	The objective function is a convex quadratic and continuous in \( \eta \).
	The constraint set is
	\begin{align*}
		\calC(\lambda)
		 & = \cbr{ \eta : \eta \in \Row(X), : \norm{\eta_\bi}_2 \leq \lambda \; \forall \; \bi \in \calB}.
	\end{align*}
	Let's show that \( \calC \) is continuous.

	Let \( \lambda_k \geq 0 \), \( \lambda_k \into \bar \lambda \)
	and \( \eta_k \in \calC(\lambda_k) \) such that \( \eta_k \into \bar \eta \).
	Since \( \eta_k \in \Row(X) \), \( \bar \eta \in \Row(X) \).
	Moreover,
	\[
		\norm{[\eta_k]_\bi}_2 \leq \lambda_k,
	\]
	so that taking
	limits on both sides implies \( \norm{\bar \eta_k}_2 \leq \bar \lambda \).
	Thus, \( \bar \eta \in \calC(\bar \lambda) \), showing that \( \calC \)
	is closed.

	\citet[Theorem 12]{hogan1973point} states that \( \calC(\lambda) \)
	is open at \( \bar \lambda \) if for each \( \bi \in \calB \), \( g_\bi(\lambda, \eta) = \norm{\eta_\bi}_2 - \lambda \)
	is continuous on \( \bar \lambda \times \calC(\lambda) \), convex in
	\( \eta \), and for fixed \( \lambda \), and there exists \( \bar \eta \)
	such that such that \( g(\lambda, \bar \eta) < 0  \).

	Let us check these conditions.
	First, observe that \( g_\bi \) is continuous and convex in \( \eta \) for
	any choice of \( \lambda \).
	Taking \( \bar \eta = 0 \), we find
	\[
		g_\bi(\lambda, 0) = -\lambda < 0,
	\]
	as long as \( \lambda > 0 \).
	Thus, \( \calC(\lambda) \) is open at each \( \lambda > 0 \).

	At \( \lambda = 0 \), \( \calC(\lambda) = \cbr{0} \).
	Since \( \calC \) is closed everywhere, we conclude it is open at
	\( \lambda = 0 \).
	Putting these results together proves that \( \calC \) is continuous.

	Recall that the dual solution satisfies \( \eta^*_\bi = \ci \),
	so that it is always unique.
	Combining this fact with \citet[Theorem 1]{robinson1974sufficient}
	implies that \( \lambda \mapsto \eta^*(\lambda) \) is a continuous
	function.
	Since we have
	\[
		\eta^*(\lambda) = c(\lambda) = X^\top (y - X w^*(\lambda)),
	\]
	it must be that the model fit \( \hat y(\lambda) = X w^*(\lambda) \)
	is continuous as well (any discontinuities must be in \( \Null(X) \),
	but \( \hat y \) is orthogonal to \( \Null(X) \)).

	Finally, using \cref{prop:value-continuity-cgl}, we have
	\[
		p^*(\lambda) - \half \norm{\hat y(\lambda) - y}_2^2 = \lambda \sum_{\bi \in \calB} \norm{\wi}_2,
	\]
	is a sum of continuous functions and thus continuous.
	Writing
	\[
		g(\lambda) = \sum_{\bi \in \calB} \norm{\wi}_2 = \sbr{\lambda \sum_{\bi \in \calB} \norm{\wi}_2} \sbr{\frac{1}{\lambda}},
	\]
	as the product of continuous functions shows
	\( g(\lambda) \) is continuous at every \( \lambda > 0 \).
\end{proof}


